% 原文 p.101（後半）


\subsectionlink{部分系への分解}
デコヒーレンスは、世界を「系」と「環境」に分けることの存在と可能性を、前提とすることに注意しよう。デコヒーレンス研究計画では、「環境」という語は、通常、系の「残り」として理解される。すなわち環境自由度は、典型的には制御しない、制御できない、または制御する必要がなく、考えている観測に直接関係しないが、それでも環境は、系の状態の発展に有意に寄与するすべての自由度を含む、という意味である。これは本質的に、Zurekのデコヒーレンスと指針状態に関する最初の論文\cite[p.~1520]{ref008}で与えられた定義である。

% 原文 p.102
\hypertarget{srcpage:102}{}\label{srcpage:102}
この系と環境の二元性は、一般に量子もつれと関係する。量子もつれは常に、宇宙の部分間の相関を記述する。宇宙を個々の部分系へ分解しない限り、測定問題はない。宇宙全体の状態ベクトル$\lvert\Psi\rangle$\footnote{この全状態を、あえて仮定するならば、である。Aulettaの反論\cite{ref138}を参照。}はSchrödinger方程式に従って決定論的に発展し、そこに解釈的な困難はない。測定問題は、宇宙の全Hilbert空間$\mathcal H$を、二空間の積$\mathcal H_1\otimes\mathcal H_2$へ分解し、一方の部分系に個別の状態を割り当てようとしたときにだけ生じる。Zurek\cite[p.~718]{ref016}は、こう表現する。
\begin{quote}
系が存在しなければ、解釈の問題は消えるように見える。系のない宇宙には、単に「収縮」が必要ない。古典的実在についての私たちの経験は、外から見た宇宙全体ではなく、その内部の系に当てはまるのである。
\end{quote}
さらに、「観測」「相関」「相互作用」といった語も、系への分割なしには、当然ほとんど意味を持たない。Zehは、いかなる観測も宇宙の一部についての無知から生じるという意味で、観測者の局所性が観測を定義し、この局所性が、量子系に起こりうる「事実」も定義すると示唆した。Landsman\cite[pp.~45--46]{ref139}も同様に論じる。
\begin{quote}
量子力学における「測定」「事実」「出来事」の本質は、当該の系のある部分を観測しないこと、あるいはその部分が無関係であることにある。（中略）無関係と宣言された、または無関係であることを強いられる部分のない世界は、事実のない世界である。
\end{quote}
しかし、測定問題の出現とデコヒーレンス研究計画の定式化に、どれほど必要に見えるとしても、宇宙を部分系へ分解するという仮定は、決して自明ではない。定義により宇宙全体は閉じた系なので、外部環境の「観測されない自由度」はない。したがって、「大域的」な準古典的観測可能量を決めるために、デコヒーレンス理論を宇宙全体へ適用することはできない。また全Hilbert空間を、どのように部分系へ分けるかについての一般的基準はない一方、系の性質と呼ばれるものの多くは、他の系との相関に依存する。この問題は、デコヒーレンスに、古典性の主観的知覚の説明を動機付けるだけでなく（たとえば8.2節、9.4節）、準古典的な「巨視的事実」の定義まで可能にさせたい場合に、とりわけ深刻になる。

% 原文 p.103
\hypertarget{srcpage:103}{}\label{srcpage:103}
Zurek\cite[p.~1820]{ref104}は、この概念的な難点を、最初期に明瞭に指摘した一人である。
\begin{quote}
特に、しばしば当然視されてきた一つの問題が、デコヒーレンス研究計画全体の基礎として、大きく立ちはだかっている。それは、古典性の出現についてのあらゆる議論で、それほど決定的な役割を果たす「系」とは何かという問いである。（中略）系とは何か、たとえば適切に大きなHilbert空間内の全体のHamiltonianが与えられたとき、それらをどう定義するかについての説得的説明は、疑いなく、きわめて有用であろう。
\end{quote}
よく提案される考えは、「絶対的」分解という概念を捨て、代わりに、異なる状態空間と、それら相対的に定義された空間間の相関を通じて現れる性質の、本質的な相対性を仮定することである（たとえば、デコヒーレンスとは関係しない提案\cite{ref140,ref141,ref142,ref143}）。系と相関についてのこの相対的な見方は、非古典的という意味で、反直観的な含意を持つ。しかし量子もつれの場合と同じく、これらの含意を、さらなる解決を要するパラドックスと見なす必要はない。自然のある性質を反直観的なものと受け入れることは、本質的に局所的な観測に基づく経験の範囲が許す限り、完全で客観的な自然の記述に至るために、実際に満足できる道である。

\subsectionlink{形式的な道具とその解釈}
ここで、デコヒーレンスの文脈でよく使われ、参照される、いくつかの形式的手順を簡単に概説する。2.15.1節では、二部分系の状態を、von Neumann測定の完了時の終状態を思わせる対角形に常に書けることを教える、Schmidt分解定理を導入する。Schmidt分解とデコヒーレンスの間の、意味のあるつながりと、見せかけのつながりの両方を論じる。2.15.2節では、密度行列のWigner表示を導入する。この表示は、位相空間でデコヒーレンスを可視化するのに、しばしば便利な道具である。続く2.15.3節では、任意に与えられた密度行列を、より大きな系の純粋状態密度行列について、トレースを取ったものと見なせることを示す。最後の2.15.4節では、演算子和形式を導入する。この形式を、より複雑なデコヒーレンス模型に適用するのは、通常実行不可能だが、環境と相互作用する系の縮約力学を記述するための、非常に一般的な枠組みを与える。
